
% !TEX root = ../main_hessenberg.tex
\section{Introduction}\label{sec:intro}

AI tools can now produce substantial Lean developments under human prompt
design~\cite{ilin2026semiautonomousvml, avigad2026mathematiciansageai}. This
shifts the practical question of how a formalization is built, and with it, the
methodological question of what a formalization is for. Lean is built on
dependent type theory, in which the Curry--Howard correspondence identifies
proofs with terms and mathematical statements with
types~\cite{deMoura2021Lean4}; in this sense, the formal development of a proof
is the proof itself. The major formalization successes of the last five
years~\cite{lean_liquid, vandoorn2023sphereeversion,
Anderson_Formalization_of_the_2023, becker2025carleson_blueprint,
bolan2025equationaltheoriesprojectadvancing, flt_lean} have been carried out
alongside reflective accounts of the practice of building such
developments~\cite{Avigad2023MathematicsAT, avigad2026mathematiciansageai}.

We report a case study that takes the Curry--Howard correspondence as its
working premise: the formal Lean development is the mathematics, and the prose
only renders it. The premise is not ours; a growing literature examines the
methodological consequences of treating formal developments as designed
mathematical artifacts~\cite{commelin2023abstractionboundariesspecdriven,
baanen2025growing, vandoorn2020maintaining, avigad2024anatomyformalproof,
ilin2026semiautonomousvml}. Our contribution is to set out the practices that
follow when the premise is enforced end to end. A network analysis of
\texttt{Mathlib}~\cite{li2026networkstructuremathlib} shows its structure to be
heavily infrastructural (\emph{e.g.}, coercions and typeclasses) rather than
strictly mathematical; the practices below apply this library-wide distinction
to a single development. We tag every declaration as mathematical content or
\emph{engineering tax}, that is, scaffolding with no mathematical counterpart.
We gather this overhead into a standalone Foundation layer that the
mathematical layer reaches only through clean interfaces. We enforce a naming
discipline under which every declaration reads aloud as a mathematical
statement. Under this arrangement the authors write no code: they steer the
formalization entirely through prompt design and mathematical judgment. The
central principle is as follows:
\begin{center}
\emph{the code must read as mathematics.}
\end{center}

As a case study, we characterize the digraphs of real-orthogonal upper
Hessenberg matrices, a small, self-contained problem that revisits a question
first considered in~\cite{severini2003a}. Nothing in the practices depends on
the problem being small; if anything, keeping engineering tax out of the
mathematical narrative matters more in a large formalization, where the
narrative is easiest to lose. Section~\ref{sec:methodology} sets out the
practices, and Section~\ref{sec:example} presents the case study, with direct
cross-references into the Lean development. The full development is available
\href{https://github.com/simoneseverini/Digraphs-of-Real-Orthogonal-Upper-Hessenberg-Matrices}{on
GitHub}; every theorem and definition below carries a link into its source.

We have also published the mathematics of Section~\ref{sec:example} on
\href{https://www.astrolabe.network/writing/hessenberg-digraphs}{\texttt{astrolabe.network}},
where every statement is cut into a card addressed by a unique hash, and the
relations between cards are themselves nodes, each carrying the semantics of
the relation it stands for. The visualization and the network analyses it
carries follow~\cite{li2026networkstructuremathlib}; the underlying design was
set out in~\cite{li2026astrolabecontentaddressablehypergraphsemantic}. Astrolabe
has since evolved substantially, and this network should be read as a prototype
instance of an evolving publication format rather than a fixed specification.


